\section{Evaluación}

\subsection{Planteamiento}
La evaluación es continua y sirve para mejorar la enseñanza y el aprendizaje \parencite[art.~20.1, p.~22]{decreto64}. Tiene tres momentos: diagnóstica (S1), formativa (cuestionarios, fichas y coevaluación del borrador) y sumativa (examen, libro y defensa). Participan el profesor, los iguales y el propio alumnado.

\subsection{Criterios de calificación}\label{sec:calificacion}
\begin{tablaSA}{colspec={X[2.3,l]Q[c,wd=30pt]X[1.7,l]X[2.4,l]}, cells={font=\footnotesize}}
  Instrumento & Peso & Criterios & Nota del instrumento (0--10) \\
  Examen escrito & 35\,\% & 1.1, 1.2, 2.1, 3.1, 5.1, 5.2, 7.2 & Suma de los puntos del examen \\
  Libro de GeoGebra y propuesta (R1) & 25\,\% & 1.1, 2.2, 6.1, 6.2, 7.2 & Media ponderada de los niveles de R1 \\
  Fichas de GeoGebra F1--F4 (R4) & 10\,\% & 1.1, 1.2, 2.1, 3.1, 3.2, 7.1, 7.2 & Media de las cuatro fichas, cada una valorada con R4 \\
  Exposición oral (R2) & 10\,\% & 1.1, 2.2, 6.1, 6.2, 8.1 & Media ponderada de los niveles de R2 \\
  Cuestionarios Q1--Q4 & 5\,\% & 1.2, 2.1, 5.2, 7.1 & Media de los cuatro: porcentaje de aciertos sobre 10, con el mejor de dos intentos \\
  Mapa conceptual (L2) & 5\,\% & 5.1 & (n.º de Sí / 3) $\times$ 10 \\
  Observación sistemática (L1) & 5\,\% & 3.2, 9.1, 9.2, 9.3 & (n.º de Sí / 6) $\times$ 10 \\
  Coevaluación y autoevaluación (R3) & 5\,\% & 7.2, 8.1, 9.1, 9.2, 9.3 & Media de las casillas auto y co, $\times$ 2,5 \\
\end{tablaSA}

La nota del libro es común al equipo, ajustada individualmente si la observación o la coevaluación muestran diferencias claras de participación.

\subsection{Retroalimentación y refuerzo}
\begin{itemize}
  \item Los cuestionarios se autocorrigen y determinan qué criterio refuerza cada estudiante en S13; el segundo intento se hace después de ese refuerzo.
  \item Las fichas se devuelven con un comentario sobre qué mejorar.
  \item El borrador del libro se revisa con la rúbrica antes de su entrega, de modo que la crítica pueda incorporarse.
  \item Tras el examen, quien no alcance el nivel adecuado en algún criterio recibe tareas específicas y una nueva oportunidad en la siguiente unidad.
\end{itemize}

\subsection{Instrumentos}\label{sec:instrumentos}

\subsubsection{R1. Rúbrica del libro de GeoGebra y la propuesta}
\begin{tablaSAlarga}{colspec={Q[l,wd=66pt]X[l]X[l]X[l]X[l]}, cells={font=\scriptsize}, colsep=3pt,
  row{1}={bg=urjc, fg=white, font=\scriptsize\bfseries}, column{1}={bg=rosa, font=\scriptsize\bfseries}, cell{1}{1}={bg=urjc}}
  Criterio (peso) & 1 · Inicial & 2 · En desarrollo & 3 · Adecuado & 4 · Excelente \\
  Modelo y método, 1.1 (30\,\%) & Sin modelo a escala o incorrecto & Modelo incompleto; método poco justificado & Modelo correcto, reutilizable con otra fecha y altura, y método explicado & Además compara uno y dos puntos de medida y valora su eficiencia \\
  Error y solución, 2.2 (20\,\%) & Una sola medida sin análisis & Repite medidas, sin análisis del error & Analiza un grado de error y elige un resultado & Además justifica la decisión según el contexto \\
  Representación, 7.2 (20\,\%) & Representaciones confusas & Claras pero sin relacionar & Integra construcción, cálculos y explicación, y cita sus fuentes & Además elige y justifica las representaciones más útiles \\
  Propuesta, 6.1 (20\,\%) & Sin cálculo de filas ni plano & Plano sin aplicar el criterio técnico & Plano y filas correctos con el criterio del IDAE & Además estima producción y plantea alternativas \\
  ODS y límites, 6.2 (10\,\%) & No se menciona & Mención general & Relaciona con los ODS y señala un límite & Analiza el impacto y varios límites del modelo \\
\end{tablaSAlarga}

\subsubsection{R2. Rúbrica de la exposición oral}
\vspace*{10pt}
\begin{tablaSAlarga}{colspec={Q[l,wd=66pt]X[l]X[l]X[l]X[l]}, cells={font=\scriptsize}, colsep=3pt,
  row{1}={bg=urjc, fg=white, font=\scriptsize\bfseries}, column{1}={bg=rosa, font=\scriptsize\bfseries}, cell{1}{1}={bg=urjc}}
  Criterio (peso) & 1 · Inicial & 2 · En desarrollo & 3 · Adecuado & 4 · Excelente \\
  Organización y rigor, 8.1 (40\,\%) & Desordenada, sin notación & Orden parcial, errores de notación & Ordenada, terminología correcta & Además adaptada al público y sin errores \\
  Método y solución, 1.1 y 2.2 (30\,\%) & No explica el método & Lo explica sin justificar & Explica y justifica el resultado & Además responde con solvencia a preguntas \\
  Propuesta, 6.1 y 6.2 (30\,\%) & No formula propuesta & Propuesta sin datos & Propuesta con datos y ODS & Además argumenta impacto y viabilidad \\
\end{tablaSAlarga}

\subsubsection{R3. Coevaluación y autoevaluación}
Escala de 1 (nunca) a 4 (siempre). La autoevaluación la completa cada estudiante en S20. En la coevaluación, un equipo asignado por el profesor valora los indicadores 7.2 y 8.1, y los compañeros del propio equipo, los indicadores 9.1, 9.2 y 9.3; si hay varias valoraciones, se hace la media. Si auto y co difieren en 2 puntos o más, el profesor lo comenta con el estudiante antes de cerrar la nota.

\begin{tablaSA}{colspec={X[l]Q[c,wd=30pt]Q[c,wd=62pt]Q[c,wd=44pt]Q[c,wd=44pt]}, cells={font=\footnotesize}}
  Indicador & Crit. & Coevalúa & Auto (1--4) & Co (1--4) \\
  Las representaciones del libro permiten entender el método & 7.2 & Otro equipo & & \\
  La explicación se sigue con orden y buena notación & 8.1 & Otro equipo & & \\
  Uso mis errores y los del equipo para revisar el razonamiento & 9.1 & Propio equipo & & \\
  Incorporo la crítica razonada para mejorar & 9.2 & Propio equipo & & \\
  Asumo mi rol y escucho el razonamiento de los demás & 9.3 & Propio equipo & & \\
\end{tablaSA}

\pista{Metacognición: ¿qué sabía antes? ¿Qué me costó más? ¿Qué haría distinto en la próxima medición?}

\subsubsection{L1. Lista de cotejo de observación y Kanban}
Salvo el indicador de rol, que tiene su propia regla, un indicador se marca Sí si se observa al menos una vez en las sesiones que la sección~\ref{sec:matrices} asigna a su criterio.

\begin{tablaSA}{colspec={X[l]Q[c,wd=30pt]Q[c,wd=30pt]Q[c,wd=30pt]}, cells={font=\footnotesize}}
  Indicador (por estudiante) & Crit. & Sí & No \\
  Investiga en GeoGebra antes de pedir ayuda & 3.2 & & \\
  Analiza un error propio y lo corrige & 9.1 & & \\
  Persevera ante una medida o construcción fallida & 9.2 & & \\
  Aporta al registro de cambios del libro la crítica recibida y lo que modifica & 9.2 & & \\
  Cumple los criterios de su rol\soloProducto{ (sección~\ref{sec:cooperativo})} en al menos tres de los cuatro turnos & 9.3 & & \\
  Escucha y tiene en cuenta el razonamiento ajeno & 9.3 & & \\
\end{tablaSA}

\subsubsection{L2. Lista de cotejo del mapa conceptual}
\begin{tablaSA}{colspec={X[l]Q[c,wd=30pt]Q[c,wd=30pt]Q[c,wd=30pt]}, cells={font=\footnotesize}}
  Indicador & Crit. & Sí & No \\
  Relaciona vector director y normal con las ecuaciones de la recta & 5.1 & & \\
  Conecta producto escalar con ángulos y distancias & 5.1 & & \\
  Indica qué herramienta usa para cada tipo de problema & 5.1 & & \\
\end{tablaSA}

\subsubsection{R4. Rúbrica común de las fichas de GeoGebra}
\begin{tablaSAlarga}{colspec={Q[l,wd=66pt]X[l]X[l]X[l]X[l]}, cells={font=\scriptsize}, colsep=3pt,
  row{1}={bg=urjc, fg=white, font=\scriptsize\bfseries}, column{1}={bg=rosa, font=\scriptsize\bfseries}, cell{1}{1}={bg=urjc}}
  Aspecto (peso) & 1 · Inicial & 2 · En desarrollo & 3 · Adecuado & 4 · Excelente \\
  Construcción (30\,\%) & No funciona & Funciona con errores o incompleta & Correcta y dinámica & Además está etiquetada y se puede reutilizar \\
  Resultado y comprobación (30\,\%) & Sin resultado & Resultado sin comprobar & Resultado correcto y comprobado a mano & Además explica las discrepancias \\
  Respuesta justificada (40\,\%) & No responde & Responde sin justificar & Justifica con un razonamiento correcto & Además generaliza o propone otro caso \\
\end{tablaSAlarga}

\pista{Criterios que recoge cada aspecto. Construcción: 7.1 en F1, 3.2 en F2 y 7.2 en F4. Resultado y comprobación: 2.1 en F1 y F3, y 7.2 en F4. Respuesta justificada: 2.1 en F1, 3.1 en F2, 1.1 y 1.2 en F3, y 7.2 en F4.}

\subsubsection{Fichas de GeoGebra}\label{sec:fichas}
\soloProducto{Las cuatro fichas siguen la progresión descrita en la sección~\ref{sec:andamiaje}. }Se valoran con la rúbrica R4.

\begin{tarjeta}[Paso a paso · sesión 2 · criterios 7.1 y 2.1]{F1. Vectores y producto escalar}
\begin{enumerate}[leftmargin=1.5em]
  \item Abre el archivo de partida, con ejes y cuadrícula. Crea $\vec u=(4,3)$ y $\vec v=(1,2)$ con la herramienta Vector.
  \item Crea $\vec u+\vec v$ y $2\vec u$. Arrastra los extremos y describe qué figura forman $\vec u$, $\vec v$ y $\vec u+\vec v$.
  \item Muestra $\vec u\cdot\vec v$ y el ángulo que forman. Arrastra $\vec v$ hasta que el producto escalar valga 0.
  \item Vuelve a $\vec v=(1,2)$. Calcula a mano $|\vec u|$, el ángulo y la proyección de $\vec u$ sobre $\vec v$, y compáralos con GeoGebra.
  \item Sea $\vec s$ el vector hacia el sol y $\vec n$ la normal al panel. ¿Cuándo es máximo $\vec s\cdot\vec n$? ¿Qué significa que valga 0?
\end{enumerate}
\campo{Pregunta final.} ¿Cuándo es nulo el producto escalar? Compruébalo a mano con un ejemplo.

\campo{Solución.} $\vec u\cdot\vec v=10$, $|\vec u|=5$, $|\vec v|=\sqrt5$, $\cos\alpha=2/\sqrt5$, $\alpha\approx26{,}57^\circ$ y proyección $10/\sqrt5\approx4{,}47$. Dos vectores no nulos tienen producto escalar nulo si son perpendiculares: $(4,3)\cdot(-3,4)=0$. El máximo se da con el sol de frente; el valor 0 corresponde a un rayo rasante.

\campo{Refuerzo:} captura de cada paso. \campo{Ampliación:} halla los vectores unitarios perpendiculares a $\vec u$.
\end{tarjeta}

\begin{tarjeta}[Guion por objetivos · sesión 5 · criterios 3.2 y 3.1]{F2. La circunferencia goniométrica}
Tu construcción debe incluir:
\begin{enumerate}[leftmargin=1.5em]
  \item Un deslizador $\alpha\in[0^\circ,360^\circ]$, la circunferencia de radio 1 y un punto $P$ que la recorra.
  \item Segmentos que muestren el seno y el coseno de $\alpha$, y el valor de la tangente.
  \item Una tabla con los signos de las tres razones en cada cuadrante.
  \item El valor de $\sin^2\alpha+\cos^2\alpha$ mientras se mueve el deslizador.
  \item El cociente $\sin 2\alpha/(\sin\alpha\cos\alpha)$. Un compañero propone que $\sin 2\alpha=2\sin\alpha$: busca un contraejemplo y propón una relación correcta.
\end{enumerate}
\campo{Pregunta final.} Conjetura una identidad y demuéstrala por escrito (la demostración se completa en S6).

\campo{Solución.} $\sin^2\alpha+\cos^2\alpha=1$: es el teorema de Pitágoras en un triángulo de hipotenusa 1, válido en todos los cuadrantes porque las razones aparecen al cuadrado. Contraejemplo: con $\alpha=30^\circ$, $\sin60^\circ\approx0{,}866\neq1$. Donde el denominador no se anula, el cociente vale 2, luego $\sin2\alpha=2\sin\alpha\cos\alpha$, que se demuestra en S7 haciendo $\beta=\alpha$ en el ángulo suma.

\campo{Refuerzo:} se da la definición $P=(\cos\alpha,\sin\alpha)$. \campo{Ampliación:} conjetura una expresión para $\cos2\alpha$.
\end{tarjeta}

\begin{tarjeta}[Encargo con pistas · sesión 9 · criterios 1.1, 1.2 y 2.1]{F3. ¿Cuántos triángulos hay?}
\campo{Encargo.} Construye un triángulo con $A=30^\circ$, $b=10$ y un deslizador para el lado $a$, de modo que se vean todas las posiciones posibles de $B$.
\begin{enumerate}[leftmargin=1.5em]
  \item ¿Para qué valores de $a$ hay dos triángulos, uno o ninguno? Justifícalo con un cálculo.
  \item Para $a=8$, resuelve a mano todos los triángulos y compruébalos arrastrando el deslizador. Si falta tiempo, este apartado se termina en casa.
  \item Si conocieras dos lados y el ángulo comprendido, ¿qué teorema usarías? ¿Y en este caso? Razónalo.
\end{enumerate}
\campo{Pistas, solo si se piden.} 1: ¿qué figura forman los puntos situados a distancia $a$ de $C$? 2: ¿a qué distancia está $C$ de la recta $AB$? 3: compara $a$ con $b\sin A$ y con $b$.

\campo{Solución.} $b\sin A=5$. Ninguno si $a<5$; uno, rectángulo, si $a=5$; dos si $5<a<10$; uno si $a\ge10$. Para $a=8$, $\sin B=0{,}625$: $B\approx38{,}68^\circ$, $C\approx111{,}32^\circ$, $c\approx14{,}91$; o bien $B\approx141{,}32^\circ$, $C\approx8{,}68^\circ$, $c\approx2{,}42$. Con dos lados y el ángulo comprendido conviene el teorema del coseno; aquí, con un lado y su ángulo opuesto, el del seno.

\campo{Refuerzo:} las pistas 1 y 2 se dan desde el inicio. \campo{Ampliación:} estudia qué ocurre si $A$ es obtuso.
\end{tarjeta}

\begin{tarjeta}[Encargo abierto · sesión 12 · criterio 7.2]{F4. Rayo y sombra}
\campo{Encargo.} Construye un modelo dinámico de la sombra de un edificio que cumpla estos requisitos:
\begin{enumerate}[leftmargin=1.5em]
  \item Deslizadores $h\in[5^\circ,85^\circ]$ y $H\in[0,20]$, y el punto $T=(0,H)$.
  \item El rayo que pasa por $T$, definido con un vector director que debes deducir a partir de $h$.
  \item El punto $S$ donde el rayo corta al eje $OX$, comprobando que su abscisa es $H/\tan h$.
  \item Un panel en el punto $P=(20,3)$, con $H=12$.
  \item La recta en forma explícita y general. ¿Cuál te sirve mejor para decidir si $P$ queda en sombra? Justifícalo.
\end{enumerate}
\campo{Pregunta final.} ¿Para qué valores de $h$ queda $P$ en sombra?

\campo{Solución.} $P$ está en sombra si $3<12-20\tan h$, es decir, si $\tan h<0{,}45$, luego $h<24{,}23^\circ$. Un vector director es $\vec v=(\cos h,-\sin h)$. La forma explícita es la más cómoda, porque permite comparar directamente la altura de la recta con la de $P$.

\campo{Tarjetas de ayuda para el equipo, si las pide:} el vector $\vec v$, los pasos con capturas y datos prefijados. \campo{Ampliación:} añade un panel inclinado y calcula $h$ para incidencia perpendicular.
\end{tarjeta}

\subsubsection{Hoja de toma de datos}
\begin{tablaSA}{colspec={X[2,l]X[1,c]X[1,c]X[1,c]X[1,c]}, cells={font=\footnotesize}}
  Magnitud (con unidad) & Medida 1 & Medida 2 & Medida 3 & Valor elegido \\
  Altura del ojo $a$ (m) & & & & \\
  Un punto de medida: distancia al pie $D$ (m) & & & & \\
  Un punto de medida: ángulo de elevación (°) & & & & \\
  Dos puntos de medida: distancia entre ellos $d$ (m) & & & & \\
  Dos puntos de medida: ángulo $\alpha$ en A (°) & & & & \\
  Dos puntos de medida: ángulo $\beta$ en B (°) & & & & \\
  Sombra de la varilla de 1~m (m), con fecha y hora & & & & \\
\end{tablaSA}

\subsubsection{Cuestionario modelo}
Banco de preguntas de autocorrección en GeoGebra Classroom, con respuesta entre paréntesis.
\begin{enumerate}[label=\textbf{\arabic*.},leftmargin=1.6em]
  \item (Q1, 7.1) Si $\vec u=(1,2)$ y $\vec v=(3,-1)$, ¿cuánto vale $\vec u\cdot\vec v$? \emph{(1)}
  \item (Q1, 5.2) ¿Para qué $k$ son perpendiculares $(2,k)$ y $(3,6)$? \emph{($k=-1$)}
  \item (Q2, 1.2) Soluciones de $2\sin x=1$ en $[0^\circ,360^\circ)$. \emph{($30^\circ$ y $150^\circ$)}
  \item (Q3, 2.1) Con $a=3$, $b=5$ y $A=80^\circ$, ¿existe el triángulo? \emph{(No: $\sin B\approx 1{,}64>1$)}
  \item (Q4, 5.2) Distancia de $(0,0)$ a $3x+4y-10=0$. \emph{(2)}
\end{enumerate}

\subsubsection{Examen modelo}\label{sec:examen}
\begin{tarjeta}[Individual · 45 minutos · calculadora permitida]{Examen. Medir lo que no se puede alcanzar}
\begin{enumerate}[label=\textbf{\arabic*.},leftmargin=1.6em,itemsep=3pt]
  \item \textbf{(2 puntos; 5.2)} Dados $\vec u=(3,1)$ y $\vec v=(2,-1)$: a) calcula $\vec u+2\vec v$ y $\vec u\cdot\vec v$; b) halla el ángulo que forman y la proyección de $\vec u$ sobre $\vec v$; c) halla $k$ para que $\vec w=(k,2)$ sea perpendicular a $\vec u$.
  \item \textbf{(1 punto; 2.1)} Un compañero afirma que la proyección de $\vec u$ sobre $\vec v$ vale $5/\sqrt{10}$, pero GeoGebra muestra $2{,}236$. ¿Quién tiene razón? Explica el origen del error.
  \item \textbf{(2 puntos; 3.1, 1.2)} a) A partir de la fórmula de $\sin(\alpha+\beta)$, demuestra que $\sin 2\alpha=2\sin\alpha\cos\alpha$. b) Resuelve $\sin 2x=\cos x$ dando todas las soluciones.
  \item \textbf{(2,5 puntos; 1.1, 1.2)} a) Desde $A$ se ve lo alto de un edificio con un ángulo de $25^\circ$ y, acercándose 10~m, con $38^\circ$. El ojo está a $1{,}6$~m del suelo. Calcula la altura y justifica qué teorema has elegido. b) Resuelve el triángulo con $a=8$, $b=10$ y $A=40^\circ$.
  \item \textbf{(2,5 puntos; 5.1, 5.2, 7.2)} a) Escribe las ecuaciones de la recta $r$ que pasa por $A(1,2)$ con vector director $(2,-1)$. b) Estudia su posición respecto a $s: 3x-4y+5=0$ y halla el ángulo que forman. c) Calcula la distancia de $P(4,3)$ a $r$. d) Un edificio de 12~m tiene su fachada norte en $x=0$ y el sol está a $26{,}15^\circ$: ¿queda en sombra un panel en $(20,3)$? Indica qué ecuación de la recta has preferido y por qué.
\end{enumerate}
\end{tarjeta}

\begin{tarjeta}{Solución del examen}
\small
\textbf{1.} a) $(7,-1)$; $\vec u\cdot\vec v=5$. b) $\cos\alpha=5/(\sqrt{10}\sqrt5)=\sqrt2/2$, luego $\alpha=45^\circ$; proyección $5/\sqrt5=\sqrt5$. c) $3k+2=0$, $k=-2/3$.

\textbf{2.} GeoGebra: $\sqrt5\approx2{,}236$. El compañero dividió por $|\vec u|$, con lo que calculó la proyección de $\vec v$ sobre $\vec u$.

\textbf{3.} a) Tomando $\beta=\alpha$. b) $\cos x(2\sin x-1)=0$: $x=90^\circ+180^\circ k$, $x=30^\circ+360^\circ k$, $x=150^\circ+360^\circ k$, con $k\in\mathbb Z$.

\textbf{4.} a) En el triángulo se conocen un lado y dos ángulos, por lo que se usa el teorema del seno: $H=\dfrac{10\sin25^\circ\sin38^\circ}{\sin13^\circ}+1{,}6\approx13{,}17$~m. b) Dos soluciones: $B\approx53{,}46^\circ$, $C\approx86{,}54^\circ$, $c\approx12{,}42$; o bien $B\approx126{,}54^\circ$, $C\approx13{,}46^\circ$, $c\approx2{,}90$.

\textbf{5.} a) $(x,y)=(1,2)+t(2,-1)$; $x=1+2t$, $y=2-t$; $\frac{x-1}{2}=\frac{y-2}{-1}$; $x+2y-5=0$; $y=-\frac12x+\frac52$. b) Secantes, pues $(1,2)$ y $(3,-4)$ no son proporcionales; $\cos\theta=\frac{|3-8|}{\sqrt5\cdot5}=\frac1{\sqrt5}$, $\theta\approx63{,}43^\circ$. c) $d=\frac{|4+6-5|}{\sqrt5}=\sqrt5$. d) En $x=20$ la recta del rayo tiene altura $12-20\tan26{,}15^\circ\approx2{,}18<3$: el panel no queda en sombra; la forma explícita permite comparar alturas directamente.
\end{tarjeta}
